% Black-hole result on dev91. Written 2026-11-16 from the receipts under
% research/flywheel_selfplay/h458_test_informed_flywheel/evidence/ (look 2 and the GEN2A look).
% Every number below carries its label. No internal representation of a solution is described here.
\subsection{Absorbing a benchmark: the dev91 reference board}
\label{sec:blackhole-dev91}

The reference board of Section~\ref{sec:results} holds 91 regression tables in 19 families, each with an official
split into training-side and test-side rows. We report three regimes and never merge them. \emph{Unseen}: the system
never trained on any board table. \emph{In-corpus}: board tables were in the training corpus of the general model,
with no per-table fitting. \emph{Absorbed}: models were fitted on all rows of the scored tables, test labels
included, always on random splits and never on the official split; at scoring time the system still reads only a
table's training-side rows. The external models saw no test row in any regime, so an absorbed result is a statement
about what a system retains of tables it has studied, not about tables it has never seen.

\begin{table}[t]
\centering
\small
\begin{tabular}{@{}llr@{}}
\hline
System & Regime & Elo \\
\hline
Stack + family models + per-table models & absorbed & 2408.5 \\
Stack + family models & absorbed & 2343.8 \\
Kumo-Tabular large & external & 2071.7 \\
TabFM ensemble & external & 1915.6 \\
Kumo-Tabular small & external & 1900.2 \\
TabFM & external & 1888.1 \\
Stack alone & in-corpus & 1859.0 \\
TabPFN 3.5 & external & 1851.6 \\
Best system that never saw the board & unseen & 1821.5 \\
\hline
\end{tabular}
\caption{The dev91 reference board under three regimes. Absorbed systems were fitted on the scored tables' rows,
test labels included, on random splits only; external models saw no test row. The top row and Kumo-Tabular large
come from one ladder fit (95\% interval of the top row: 2237 to 2726); the other rows come from the ladder fit of
the second row, in which Kumo-Tabular large is 2070.2. Intervals are about 150 to 250 Elo wide.}
\label{tab:blackhole-dev91}
\end{table}

\paragraph{Result.} With one fitted model per table family the absorbed system reaches Elo 2343.8, first on the
ladder. Adding one fitted model per table reaches 2408.5. The second step was registered before any training and
passed its paired test: all 19 families gain, by 0.38 in mean utility (a 32\% lower error on average), with an
exact sign-flip $p = 1.9\times10^{-6}$. Head to head, the absorbed system beats TabPFN 3.5 on 18 of 19 families,
TabFM on 17, Kumo-Tabular small on 19 and Kumo-Tabular large on 16.

\paragraph{What the other regimes say.} The same stack without fitted models scores 1859.0 in-corpus: above
TabPFN 3.5 on the ladder by 7 Elo, but it loses to it head to head on 15 of 19 families, and it is below TabFM and
both Kumo models. Our best system that never saw the board scores 1821.5, below all three. The gain therefore comes
from absorption, and we claim nothing here about unseen tables.

\paragraph{Splits.} For the family models, an exact audit of the training ledgers found that none of 61,814
training episodes used an official split, on any of the 91 tables. The same audit for the per-table models is
pending; their logs show no repeated split within a run.

\paragraph{Cost.} A per-table model takes 4 to 16 minutes to fit on one 24\,GB card; the 81 new ones took 3.2 hours
on four card slots. Scoring a table takes 28\,s end to end on average, of which 0.56\,s is on the GPU: the reader
processes 3.9 million table cells per GPU-second, and feature preparation on the CPU is the rest.
